% GENERATED FILE. Edit canonical lessons, then run npm run generate:books.
% canonical-source-hash: fnv1a64:401160a6edb4f735
% canonical-lessons: GE-C404-reifen, GE-C404-platten, GE-C404-aufpumpen, GE-C404-wagenheber, GE-C404-warndreieck

\chapter{Ein platter Reifen}
\label{ch:ge-ein-platter-reifen}
\hlchaptermodality{404}

\begin{chapteropening}
\textbf{By the end of this chapter:} \emph{I can deal with a flat tyre: the tyre, the flat, pumping it up, the jack and the warning triangle.}

Complete the last lesson of chapter 404: I can deal with a flat tyre: the tyre, the flat, pumping it up, the jack and the warning triangle.
\end{chapteropening}

\section[der Reifen]{der Reifen — the tyre (US: tire)}

\label{lesson:GE-C404-reifen}



\begin{quote}
Before the new one: say the German for the tow truck, then the German for the mechanic.
\end{quote}

\subsection*{You'll want to know: der Reifen}
\textbf{der Reifen} — \textquotedblleft{}the tyre (US: tire)\textquotedblright{}.

\begin{quote}
Im Herbst wechseln wir die Reifen.
\end{quote}

\begin{quote}
\emph{In autumn we change the tyres.}
\end{quote}

Same form in the plural. \emph{Die Winterreifen} — winter tyres.

\begin{cousinweb}
\emph{Reifen} goes back to \emph{Reif} 'ring, hoop', inherited from Old High German \emph{reif} 'band, ring' (EtymDB). EtymDB links it with English \emph{rope}, from the same Germanic word for a band.
\end{cousinweb}

\begin{grammarlens}[title={What you've built}]
One more word for telling, explaining and describing at B1.
\end{grammarlens}

\subsection*{Your turn}
\begin{itemize}
\raggedright
  \item \emph{Say it:} \emph{der Reifen}
  \item \emph{Say it:} \emph{der Reifen}, once more
  \item \emph{Say it:} a sentence of your own with \emph{der Reifen}
  \item \emph{Recall:} say the German for the tow truck, then the German for the mechanic, then say \emph{der Reifen} again
\end{itemize}

\begin{tcolorbox}[breakable,colback=teal!4,colframe=teal!35!black,title={Before you move on}]
What does der Reifen mean? (\textquotedblleft{}The tyre (US: tire)\textquotedblright{}.) Say it once more.
\end{tcolorbox}
\section[der Platten]{der Platten — the flat tyre (colloquial)}

\label{lesson:GE-C404-platten}



\begin{quote}
Before the new one: say the German for the mechanic, then the German for the tyre.
\end{quote}

\subsection*{You'll want to know: der Platten}
\textbf{der Platten} — \textquotedblleft{}the flat tyre (colloquial)\textquotedblright{}.

\begin{quote}
Auf dem Weg zur Arbeit hatte sie einen Platten.
\end{quote}

\begin{quote}
\emph{On the way to work she had a flat tyre.}
\end{quote}

Everyday speech. More formally: \emph{ein platter Reifen}.

\begin{cousinweb}
\emph{Der Platten} is a short everyday form of \emph{ein platter Reifen}, built on \emph{platt} 'flat'. EtymDB links \emph{platt} with French \emph{plat} and Vulgar Latin \emph{plattus} 'flat'.
\end{cousinweb}

\begin{grammarlens}[title={What you've built}]
One more word for telling, explaining and describing at B1.
\end{grammarlens}

\subsection*{Your turn}
\begin{itemize}
\raggedright
  \item \emph{Say it:} \emph{der Platten}
  \item \emph{Say it:} \emph{der Platten}, once more
  \item \emph{Say it:} a sentence of your own with \emph{der Platten}
  \item \emph{Recall:} say the German for the mechanic, then the German for the tyre, then say \emph{der Platten} again
\end{itemize}

\begin{tcolorbox}[breakable,colback=teal!4,colframe=teal!35!black,title={Before you move on}]
What does der Platten mean? (\textquotedblleft{}The flat tyre (colloquial)\textquotedblright{}.) Say it once more.
\end{tcolorbox}
\section[aufpumpen]{aufpumpen — to pump up}

\label{lesson:GE-C404-aufpumpen}



\begin{quote}
Before the new one: say the German for the tyre, then the German for the flat tyre.
\end{quote}

\subsection*{You'll want to know: aufpumpen}
\textbf{aufpumpen} — \textquotedblleft{}to pump up\textquotedblright{}.

\begin{quote}
Er pumpt die Reifen von seinem Fahrrad auf.
\end{quote}

\begin{quote}
\emph{He pumps up the tyres on his bike.}
\end{quote}

Separable: \emph{ich pumpe auf}, \emph{ich habe aufgepumpt}. \emph{Die Luftpumpe} — the pump.

\begin{cousinweb}
\emph{aufpumpen} is built from \emph{auf} 'up' and \emph{pumpen}, from \emph{Pumpe} 'pump', a seafarers' word that EtymDB traces to Middle Low German \emph{pumpe}; its deeper origin is not settled.
\end{cousinweb}

\begin{grammarlens}[title={What you've built}]
One more word for telling, explaining and describing at B1.
\end{grammarlens}

\subsection*{Your turn}
\begin{itemize}
\raggedright
  \item \emph{Say it:} \emph{aufpumpen}
  \item \emph{Say it:} \emph{aufpumpen}, once more
  \item \emph{Say it:} a sentence of your own with \emph{aufpumpen}
  \item \emph{Recall:} say the German for the tyre, then the German for the flat tyre, then say \emph{aufpumpen} again
\end{itemize}

\begin{tcolorbox}[breakable,colback=teal!4,colframe=teal!35!black,title={Before you move on}]
What does aufpumpen mean? (\textquotedblleft{}To pump up\textquotedblright{}.) Say it once more.
\end{tcolorbox}
\section[der Wagenheber]{der Wagenheber — the jack (for a car)}

\label{lesson:GE-C404-wagenheber}



\begin{quote}
Before the new one: say the German for the flat tyre, then the German for to pump up.
\end{quote}

\subsection*{You'll want to know: der Wagenheber}
\textbf{der Wagenheber} — \textquotedblleft{}the jack (for a car)\textquotedblright{}.

\begin{quote}
Mit dem Wagenheber hebt sie das Auto ein bisschen hoch.
\end{quote}

\begin{quote}
\emph{With the jack she lifts the car up a little.}
\end{quote}

Same form in the plural.

\begin{cousinweb}
\emph{Wagenheber} is built from \emph{Wagen} 'car, vehicle' and \emph{Heber} 'lifter', from \emph{heben} 'to lift'. Both are inherited: \emph{Wagen} from Proto-Germanic \emph{wagnaz} and \emph{heben} from \emph{habjaną} (EtymDB).
\end{cousinweb}

\begin{grammarlens}[title={What you've built}]
One more word for telling, explaining and describing at B1.
\end{grammarlens}

\subsection*{Your turn}
\begin{itemize}
\raggedright
  \item \emph{Say it:} \emph{der Wagenheber}
  \item \emph{Say it:} \emph{der Wagenheber}, once more
  \item \emph{Say it:} a sentence of your own with \emph{der Wagenheber}
  \item \emph{Recall:} say the German for the flat tyre, then the German for to pump up, then say \emph{der Wagenheber} again
\end{itemize}

\begin{tcolorbox}[breakable,colback=teal!4,colframe=teal!35!black,title={Before you move on}]
What does der Wagenheber mean? (\textquotedblleft{}The jack (for a car)\textquotedblright{}.) Say it once more.
\end{tcolorbox}
\section[das Warndreieck]{das Warndreieck — the warning triangle}

\label{lesson:GE-C404-warndreieck}



\begin{quote}
Before the new one: say the German for to pump up, then the German for the jack.
\end{quote}

\subsection*{You'll want to know: das Warndreieck}
\textbf{das Warndreieck} — \textquotedblleft{}the warning triangle\textquotedblright{}.

\begin{quote}
Nach der Panne stellt er das Warndreieck hinter das Auto.
\end{quote}

\begin{quote}
\emph{After the breakdown he puts the warning triangle behind the car.}
\end{quote}

Plural: \emph{die Warndreiecke}. By law every car carries one.

\begin{cousinweb}
\emph{Warndreieck} is built from \emph{warnen} 'to warn', inherited from Middle High German \emph{warnen} and a cousin of English \emph{warn} (EtymDB), and \emph{Dreieck} 'triangle', from \emph{drei} 'three' and \emph{Eck} 'corner'.
\end{cousinweb}

\begin{grammarlens}[title={What you've built}]
One more word for telling, explaining and describing at B1.
\end{grammarlens}

\subsection*{Your turn}
\begin{itemize}
\raggedright
  \item \emph{Say it:} \emph{das Warndreieck}
  \item \emph{Say it:} \emph{das Warndreieck}, once more
  \item \emph{Say it:} a sentence of your own with \emph{das Warndreieck}
  \item \emph{Recall:} say the German for to pump up, then the German for the jack, then say \emph{das Warndreieck} again
\end{itemize}

\begin{tcolorbox}[breakable,colback=teal!4,colframe=teal!35!black,title={Before you move on}]
What does das Warndreieck mean? (\textquotedblleft{}The warning triangle\textquotedblright{}.) Say it once more.
\end{tcolorbox}
